\section*{Capítulo \ref{ch:g12:exp} --- Exponencial y logaritmo}

\begin{solution}{exo:g12:exp:1}
$\dfrac{\eu^{3x}\,\eu^{-x+1}}{\eu^{x}} = \eu^{3x - x + 1 - x} = \eu^{x+1}$.

$\ln\!\left(\dfrac{\eu^2\sqrt{\eu}}{\eu^{-3}}\right)
= 2 + \tfrac12 + 3 = \tfrac{11}{2}$.
\end{solution}

\begin{solution}{exo:g12:exp:2}
\emph{(a)} Tomamos $X = \eu^x > 0$: $X^2 - 3X + 2 = 0$ da $X = 1$ o
$X = 2$, los dos positivos, luego $x = 0$ o $x = \ln 2$.

\emph{(b)} \omterm{def:g11:func:function}{Dominio}: $x - 1 > 0$ y $x + 2 > 0$, luego $x > 1$. La \omterm{def:g10:algebra:equation}{ecuación}
queda $\ln\bigl((x-1)(x+2)\bigr) = \ln 4$, de donde $(x-1)(x+2) = 4$, es
decir, $x^2 + x - 6 = 0$, luego $x = 2$ o $x = -3$. Solo $x = 2$ está en el
\omterm{def:g11:func:function}{dominio}: $S = \{2\}$.
\end{solution}

\begin{solution}{exo:g12:exp:3}
Dividimos entre $\eu^x$:
$\dfrac{1 - x^2\eu^{-x}}{1 + \eu^{-x}} \to \dfrac{1-0}{1+0} = 1$, usando
$x^2 \eu^{-x} \to 0$ (\cref{thm:g12:exp:growth}).

$\dfrac{\ln x}{\sqrt x} \to 0$ por la \cref{prop:g12:exp:lnanalytic} (caso
$n = 2$).

$x^2 \ln x = x \cdot (x \ln x) \to 0 \times 0 = 0$.

$x - \ln x = x\left(1 - \frac{\ln x}{x}\right) \to +\infty$, porque
$\frac{\ln x}{x} \to 0$.
\end{solution}

\begin{solution}{exo:g12:exp:4}
$f'(x) = \eu^{-x} - x\eu^{-x} = (1 - x)\eu^{-x}$, con el signo de $1 - x$:
$f$ crece en $\intoc{-\infty}{1}$ y decrece en $\intco{1}{+\infty}$, con un
\omterm{def:g10:functions:extrema}{máximo} absoluto $f(1) = \eu^{-1}$.

Límites: cuando $x \to +\infty$, $f(x) = \frac{x}{\eu^x} \to 0$
(\cref{thm:g12:exp:growth}); cuando $x \to -\infty$, $x \to -\infty$ y
$\eu^{-x} \to +\infty$, luego $f(x) \to -\infty$.

$f''(x) = -\eu^{-x} - (1-x)\eu^{-x} = (x - 2)\eu^{-x}$, que cambia de signo
en $x = 2$: \omterm{def:g12:deriv:inflection}{punto de inflexión} en $\bigl(2,\, 2\eu^{-2}\bigr)$.
\end{solution}

\begin{solution}{exo:g12:exp:5}
Sea $g(x) = x - \ln(1+x)$ en $\intoo{-1}{+\infty}$. Entonces
$g'(x) = 1 - \frac{1}{1+x} = \frac{x}{1+x}$, negativa en $\intoo{-1}{0}$ y
positiva en $\intoo{0}{+\infty}$: $g$ alcanza su \omterm{def:g10:functions:extrema}{mínimo} $g(0) = 0$, luego
$g \geq 0$ con igualdad solo en $0$. Por tanto, $\ln(1+x) \leq x$.

Aplicándolo con $x = \frac1n$: $n\ln\left(1 + \frac1n\right) \leq 1$, luego
$\left(1+\frac1n\right)^n = \eu^{\,n\ln(1+1/n)} \leq \eu$. Y aplicándolo
con $x = -\frac{1}{n+1} > -1$:
$\ln\left(1 - \frac{1}{n+1}\right) \leq -\frac{1}{n+1}$, luego
$-(n+1)\ln\left(1 - \frac{1}{n+1}\right) \geq 1$ y
$\left(1 - \frac{1}{n+1}\right)^{-(n+1)} \geq \eu$.
\end{solution}

\begin{solution}{exo:g12:exp:6}
\emph{1.} Cada año multiplica el capital por $1.03$: $C_n = C_0 (1.03)^n$.

\emph{2.} $C_n \geq 2C_0 \iff (1.03)^n \geq 2 \iff n \ln 1.03 \geq \ln 2
\iff n \geq \frac{\ln 2}{\ln 1.03} \approx 23.45$: el capital se duplica al
cabo de $24$ años.

\emph{3.} $\frac{70}{3} \approx 23.3$, cerca del valor exacto
$\frac{\ln 2}{\ln 1.03}$. La regla funciona porque
$\ln(1 + r) \approx r$ para $r$ pequeño, así que
$\frac{\ln 2}{\ln(1+r)} \approx \frac{0.693}{r} \approx \frac{70}{100r}$.
\end{solution}

\begin{solution}{exo:g12:exp:7}
$\eu^{2x} > \eu^{x+1} \iff 2x > x + 1 \iff x > 1$ (la exponencial es
estrictamente \omterm{def:g12:seq:monotonic}{creciente}): $S = \intoo{1}{+\infty}$.

\omterm{def:g11:func:function}{Dominio} de la segunda inecuación: $x^2 - 1 > 0$ y $x + 5 > 0$, es decir,
$x \in \intoo{-5}{-1} \cup \intoo{1}{+\infty}$. En ese \omterm{def:g11:func:function}{dominio}, como $\ln$
es estrictamente \omterm{def:g12:seq:monotonic}{creciente},
\[
\ln(x^2-1) \leq \ln(x+5) \iff x^2 - 1 \leq x + 5 \iff x^2 - x - 6 \leq 0
\iff x \in \intcc{-2}{3}.
\]
Intersecando con el \omterm{def:g11:func:function}{dominio}: $S = \intco{-2}{-1} \cup \intoc{1}{3}$.
\end{solution}

\begin{solution}{exo:g12:exp:8}
\emph{1.} $f'(x) = \dfrac{\eu^x x - \eu^x}{x^2} = \dfrac{(x-1)\eu^x}{x^2}$,
negativa en $\intoo{0}{1}$ y positiva en $\intoo{1}{+\infty}$: \omterm{def:g10:functions:extrema}{mínimo}
$f(1) = \eu$. Límites: $f \to +\infty$ en $0^+$ (el numerador $\to 1$ y el
denominador $\to 0^+$) y en $+\infty$ (\cref{thm:g12:exp:growth}).

\emph{2.} Para $x > 0$, $\eu^x = kx \iff f(x) = k$. Por la \omterm{def:g10:functions:table}{tabla de
variación} ($+\infty \searrow \eu \nearrow +\infty$, \omterm{def:g12:limcont:continuity}{continua} y
estrictamente \omterm{def:g12:seq:monotonic}{monótona} a cada lado): ninguna solución para $k < \eu$;
exactamente una ($x = 1$) para $k = \eu$; y exactamente dos para $k > \eu$
(una en $\intoo{0}{1}$ y otra en $\intoo{1}{+\infty}$, por el teorema de la
biyección en cada \omterm{def:g10:numbers:interval}{intervalo}).
\end{solution}

\begin{solution}{exo:g12:exp:9}
\emph{1.} Directamente de la primera desigualdad del
\cref{exo:g12:exp:5}: $u_n = \eu^{\,n\ln(1+1/n)} \leq \eu^1 = \eu$.

\emph{2.} Escribimos
\[
\ln u_n = n \ln\!\left(1 + \frac1n\right)
= \frac{\ln\!\left(1 + \frac1n\right) - \ln 1}{\frac1n},
\]
que es un cociente incremental de $\ln$ en el punto $1$ con incremento
$h = \frac1n \to 0$. Como $\ln$ es \omterm{def:g12:deriv:derivative}{derivable} en $1$ con \omterm{def:g12:deriv:derivative}{derivada} $1$,
$\ln u_n \to 1$. Por la \omterm{def:g12:limcont:continuity}{continuidad} de $\exp$
(\cref{prop:g12:limcont:seqcont}), $u_n = \eu^{\ln u_n} \to \eu^1 = \eu$.
\end{solution}

\begin{solution}{pb:g12:exp:1}
\textbf{1.} $x = \frac{\ln 5}{2}$. \quad $3x - 1 = \eu^2$:
$x = \frac{\eu^2 + 1}{3}$. \quad Con $u = \eu^x$:
$u^2 - 3u + 2 = (u - 1)(u - 2) = 0$: $x = 0$ o $x = \ln 2$.

\textbf{2.} $\frac{3\ln 2}{\ln 2} = 3$; \quad $3 + \frac12 = \frac72$;
\quad $\frac52$.

\textbf{3.} $f'(x) = \eu^x - 1$: negativa antes de $0$ y positiva después:
\omterm{def:g10:functions:extrema}{mínimo} $f(0) = 0$, luego $f \geq 0$ en todas partes: $\eu^x \geq 1 + x$.
Sustituyendo $x = \ln(1 + u)$: $1 + u \geq 1 + \ln(1 + u)$, es decir,
$\ln(1 + u) \leq u$.

\textbf{4.} $x^2\eu^{-x} = \frac{x^2}{\eu^x} \to 0$ y
$\frac{\ln x}{x} \to 0$: las exponenciales aplastan a las potencias y las
potencias, a los \omterm{def:g12:exp:ln}{logaritmos}. Y
$\frac{\eu^x - 1}{x} = \frac{\eu^x - \eu^0}{x - 0} \to \exp'(0) = 1$.

\textbf{5.} $f'(x) = \ln x + 1$: se anula en $\eu^{-1}$, es negativa antes
y positiva después: \omterm{def:g10:functions:extrema}{mínimo} $f\!\left(\frac1\eu\right) = -\frac1\eu$; y
$x \ln x \to 0$ en $0^+$: la curva sale del origen, baja hasta
$-\frac1\eu$ y sube después.

\textbf{6.} $n = \frac{\ln 2}{\ln 1.03} \approx 23.4$ años.

\textbf{7.} Exactos: $23.4$, $11.9$ y $8.0$ años; regla del 72: $24$, $12$
y $8$. Como $\ln(1+r) \approx r$,
$\frac{\ln 2}{\ln(1+r)} \approx \frac{0.693}{r}$, es decir,
$\frac{69.3}{\text{tipo en }\%}$; los banqueros redondean al $72$ porque se
divide estupendamente entre $2, 3, 4, 6, 8, 9$ y $12$: el cálculo mental
gana a un decimal de precisión.

\textbf{8.} $2^{t/20} = \eu^{t \ln 2 / 20}$: $\lambda = \frac{\ln 2}{20}$
por minuto. Al cabo de $24 \times 60 = 1440$ minutos:
$2^{72} \approx 4.7 \times 10^{21}$ bacterias, más que los granos de arena
de la Tierra, a partir de una sola célula y en un día. El modelo es honrado
solo mientras duren el alimento y el espacio: el crecimiento real se dobla
en la curva logística en forma de S (la curva del epidemiólogo del
\cref{pb:g12:deriv:1}).

\textbf{9.} A las 23:00 (a las $8$ horas):
$100 \times 2^{-8/5} \approx 33$ mg. Por debajo de $10$ mg:
$2^{-t/5} < 0.1$ da $t > 5\,\frac{\ln 10}{\ln 2} \approx 16.6$ h: hacia las
7:40 de la mañana siguiente; el café de las tres tiene una cola larga.

\textbf{10.} Anual: $2$. Mensual:
$\left(1 + \frac{1}{12}\right)^{12} \approx 2.613$. Diaria:
$\approx 2.715$. El techo: $\eu = 2.71828\ldots$ (\cref{exo:g12:exp:9}):
capitalizando de forma \omterm{def:g12:limcont:continuity}{continua}, la codicia \omterm{def:g12:seq:limit}{converge}.

\textbf{11.} Si los casos crecen exponencialmente,
$c(t) = c_0 \eu^{\lambda t}$, entonces
$\ln c(t) = \ln c_0 + \lambda t$: una recta de \omterm{def:g10:reffunc:affine}{pendiente} $\lambda$, la tasa
de crecimiento. Un tramo recto en la representación logarítmica \emph{es}
la fase exponencial, y su inclinación es el tempo de la epidemia.

\textbf{12.} $120 - 60 = 60$ dB significa
$10\log_{10}(I_2/I_1) = 60$: el concierto es $10^6$ (un millón de veces)
más intenso que la conversación.

\textbf{13.} Amplitud: $10^{7-5} = 100$. Energía: $100^{3/2} = 1000$: dos
puntos de magnitud esconden tres órdenes de magnitud de energía.

\textbf{14.} $10^{6.5 - 2} = 10^{4.5} \approx 32\,000$: el zumo de limón es
treinta mil veces más ácido que la leche; el pH comprime los abismos de la
química en una escala de bolsillo.

\textbf{15.} $\log_2\!\frac{4186}{27.5} = \frac{\ln(4186/27.5)}{\ln 2}
\approx 7.25$: un piano abarca algo más de siete octavas. Los sentidos
responden a \emph{razones} de estímulo (duplicar la intensidad se siente
como un escalón, sea cual sea el nivel de partida), así que la percepción
está construida sobre una escala logarítmica, y también lo están las
unidades que inventamos para ella.

\textbf{16.} Colocar la regla de $a$ a continuación de la posición de $b$
suma las longitudes $\ln a + \ln b$, y la graduación situada en esa
longitud total marca $\eu^{\ln a + \ln b} = ab$: la regla de cálculo
calcula productos sumando \omterm{def:g12:exp:ln}{logaritmos};
$\ln(ab) = \ln a + \ln b$ (\cref{prop:g12:exp:lnalg}) tallado en madera.

\textbf{17.} El \omterm{def:g10:functions:extrema}{mínimo} de $\frac{\eu^x}{x}$ en $\intoo{0}{+\infty}$ es
$\eu$, en $x = 1$: ninguna solución para $k < \eu$, exactamente una para
$k = \eu$ y dos para $k > \eu$. Comprobación con la \omterm{def:g12:deriv:tangent}{tangente}: en $a = 1$,
la \omterm{def:g12:deriv:tangent}{tangente} a $\eu^x$ es $y = \eu + \eu(x - 1) = \eu x$, que pasa por el
origen; es la recta umbral, que roza la curva en $(1, \eu)$.

\textbf{18.} $n \ln\!\left(1 - \frac1n\right) =
\frac{\ln(1 - 1/n)}{-1/n} \times (-1) \to -1$, luego
$\left(1 - \frac1n\right)^n \to \eu^{-1} \approx 0.368$: el
$0.999^{1000} \approx 0.368$ del \cref{pb:g11:binom:1}, explicado; la
constante del ``por poco'' de la lotería es $\frac1\eu$.

\textbf{19.} Las dos son la forma límite de ``muchos \omterm{def:g11:prob:model}{sucesos}
independientes, cada uno de ellos improbable'': la \omterm{def:g10:proba:distribution}{probabilidad} de que
\emph{ninguna} de $n$ opciones de tamaño $\frac1n$ se dispare tiende a
$\eu^{-1}$, y la regla de la secretaria ajusta su ventana de rechazo para
que el éxito se concentre en esa misma constante.
$\frac1\eu \approx 0.368$.

\textbf{20.} El techo de la capitalización; la única base cuya \omterm{def:g12:deriv:tangent}{tangente} en
$0$ tiene \omterm{def:g10:reffunc:affine}{pendiente} $1$ (por eso el cálculo la adora); el crecimiento que
deja atrás a todas las potencias; la constante en la que se asientan los
``por poco'' y la parada óptima; y el \omterm{def:g12:exp:ln}{logaritmo}: la multiplicación
convertida en suma y los rangos más salvajes del mundo doblados sobre una
regla.
\end{solution}
